\section{Crystal Visualization Geometry}

PyTex treats crystal visualization as a rendering layer over the canonical crystallographic model, not as a second private geometry system.

\subsection{Atomic Coordinates}

Atomic sites are stored in fractional coordinates relative to the unit cell. For a direct basis matrix
\[
\mathbf{A} = [\mathbf{a}\ \mathbf{b}\ \mathbf{c}],
\]
a site at fractional coordinates $\mathbf{f}$ is rendered at Cartesian crystal-space position
\[
\mathbf{r} = \mathbf{A}\mathbf{f}.
\]

Supercell repetition uses explicit integer translations in the direct-lattice basis:
\[
\mathbf{r}_{ijk} = \mathbf{A}(\mathbf{f} + [i, j, k]^T).
\]

\subsection{Lattice, Cell Overlays, Planes, And Directions}

The current viewer draws lattice edges from the repeated supercell box. Optional unit-cell overlays are generated directly from translated copies of the direct-basis parallelepiped. For a cell-corner anchor $\mathbf{f}_0$ and integer span vector $\mathbf{s}$, the overlay vertices are
\[
\mathbf{r}_{\mathrm{cell}} = \mathbf{A}(\mathbf{f}_0 + \boldsymbol{\delta}),
\]
where $\boldsymbol{\delta}$ ranges over the eight parallelepiped corners implied by $\mathbf{s}$.

For hexagonal-axis lattices, PyTex can also render an auxiliary hexagonal prism for teaching and visual interpretation. Let the direct basis be
\[
\mathbf{A} = [\mathbf{a}\ \mathbf{b}\ \mathbf{c}]
\]
with $|\mathbf{a}| = |\mathbf{b}|$, $\alpha = \beta = 90^{\circ}$, and $\gamma = 120^{\circ}$. From a lattice-point anchor $\mathbf{r}_0$, the basal hexagon is formed from the six vertices
\[
\mathbf{r}_0 \pm \mathbf{a},\qquad
\mathbf{r}_0 \pm \mathbf{b},\qquad
\mathbf{r}_0 \pm (\mathbf{a}+\mathbf{b}),
\]
ordered cyclically in the basal plane and extruded by integer multiples of $\mathbf{c}$.

This hexagonal prism is an explicit visualization overlay, not a replacement for the canonical cell semantics of the PyTex data model. It is useful because it exposes the sixfold basal symmetry more directly than the direct-basis parallelepiped, but users should not confuse it with the primitive cell stored by the library.

Plane overlays are defined from Miller indices $(hkl)$ and the reciprocal-space normal
\[
\mathbf{n}_{hkl} = h \mathbf{a}^{*} + k \mathbf{b}^{*} + l \mathbf{c}^{*}.
\]

The viewer computes the polygon formed by intersecting that plane with the repeated cell bounding box. This yields a bounded geometric overlay suitable for teaching and publication figures while preserving the exact lattice and reciprocal conventions from the core model.

Direction overlays are defined from a crystallographic direction vector and an explicit fractional anchor point inside the repeated cell. If the direct-basis fractional ray is
\[
\mathbf{f}(t) = \mathbf{f}_0 + t\mathbf{u},
\]
PyTex clips the ray to the repeated-cell volume and renders the bounded segment in Cartesian crystal coordinates. This keeps direction graphics tied to the same lattice semantics as the structure itself.

\subsection{Scientific Labels}

Plane and direction annotations use a shared Miller-notation formatter. Negative indices are rendered with bar notation, for example
\[
[11\bar{2}0], \qquad (11\bar{2}1),
\]
so graphics, documentation, and future figure exporters can share one scientific text convention rather than encoding ad hoc minus-sign styles per subsystem.

\subsection{View Control}

Camera angles are rendering choices, not crystallographic semantics. When PyTex accepts a `CrystalDirection` as a view target, the direction is first interpreted in crystal coordinates and only then converted into visualization-camera angles. This keeps the distinction between scientific direction meaning and graphics control explicit.

\subsection{Current Limits}

\begin{itemize}
  \item The bond model is heuristic and based on covalent-radius proximity rather than exhaustive chemistry rules.
  \item The current backend is a static Matplotlib 3D renderer.
  \item Slab views are currently implemented as deterministic geometric filtering rather than a full clipping-volume engine.
\end{itemize}
